\capitulo{6}{Trabajos relacionados}

Este capítulo relaciona el proyecto desarrollado con otras herramientas que también trabajan con modelos de inteligencia artificial y flujos de procesamiento de imágenes. El objetivo no es hacer un estudio exhaustivo de todas las soluciones existentes, sino mostrar de forma sencilla en qué se parece y en qué se diferencia el TFG respecto a algunas plataformas representativas.


\section{Herramientas y plataformas relacionadas}

\textbf{Gradio} es una biblioteca de Python orientada a crear interfaces \textit{web} para modelos de aprendizaje automático de forma rápida \cite{gradio_docs}. Su finalidad principal es facilitar demos interactivas, con entradas y salidas visuales, sin construir una aplicación \textit{web} completa desde cero.

\textbf{Roboflow} es una plataforma centrada en visión por computador. Ofrece herramientas para preparar conjuntos de datos, entrenar modelos, desplegarlos y construir \textit{workflows} visuales mediante varios pasos de procesamiento \cite{roboflow_workflows_docs}. Es la referencia más cercana al proyecto por la idea de combinar operaciones de visión por computador en un flujo.

\textbf{Kubeflow Pipelines} es una plataforma para definir y ejecutar flujos de aprendizaje automático formados por componentes encadenados \cite{kubeflow_pipelines_docs}. Está orientada a entornos de mayor escala y a infraestructuras basadas en Kubernetes, por lo que su objetivo es más amplio que el de una aplicación \textit{web} de análisis de imágenes.

\section{Comparación con el proyecto desarrollado}

La Tabla~\ref{tabla:trabajos-relacionados} resume las principales diferencias entre las soluciones revisadas y el TFG desarrollado.

\begin{table}[htbp]
	\centering
	\small
	\begin{tabular}{p{0.22\textwidth} p{0.31\textwidth} p{0.35\textwidth}}
		\toprule
		\textbf{Solución} & \textbf{Enfoque principal} & \textbf{Diferencia con el TFG} \\
		\midrule
		Gradio & Creación rápida de interfaces \textit{web} para modelos. & El TFG no se limita a una demo; incorpora usuarios, permisos, persistencia, administración y resultados temporales. \\
		\midrule
		Roboflow & Plataforma de visión por computador con \textit{workflows}. & El TFG implementa una solución propia, con \textit{pipelines} almacenados en base de datos y gestionables desde administración. \\
		\midrule
		Kubeflow Pipelines & Orquestación de flujos de aprendizaje automático. & Kubeflow está orientado a Kubernetes; el TFG aplica la idea de ejecución por etapas a una plataforma específica de análisis de imágenes\\
		\bottomrule
	\end{tabular}
	\caption[Comparación con trabajos relacionados.]{Comparación entre las soluciones relacionadas y el proyecto desarrollado.}
	\label{tabla:trabajos-relacionados}
\end{table}

Frente a Gradio, aporta una aplicación más completa, con autenticación, persistencia, permisos y administración. Frente a Roboflow, desarrolla una solución propia centrada en \textit{pipelines} configurables dentro de la propia aplicación. Frente a Kubeflow Pipelines, comparte la idea de ejecución por etapas, pero la orienta a una plataforma \textit{web} específica para análisis de imágenes.

